% !TEX root = ../elteikthesis_en.tex
\chapter{Software Requirements}
\label{ch:requirements}

This chapter defines what the ELTE IK Assistant must do and under what constraints it
must operate. The functional requirements describe the capabilities the system exposes
to the user. The non-functional requirements define the resource, privacy, latency, and
reliability constraints that any correct implementation must satisfy. Interaction logging
and the offline data pipeline are internal system concerns and do not appear as
user-facing requirements.

% ─────────────────────────────────────────────────────────────
\section{Functional Requirements}
\label{sec:req-functional}

\begin{figure}[htbp]
    \centering
    \includegraphics[width=0.8\textwidth]{UseCaseDiagram.png}
    \caption{Use case diagram for the ELTE IK Assistant}
    \label{fig:use-case}
\end{figure}

Figure~\ref{fig:use-case} illustrates the primary interactions between the user and the
system. The single actor is the \textbf{User}, representing any student or staff member
accessing the assistant through the browser interface. Asking a question is the central
use case: it always includes receiving a grounded answer with source references, and
optionally extends to selecting a language model before submitting. Document upload
always includes indexing the uploaded file into the vector store. Viewing session history
optionally extends to continuing or deleting an existing session.

\begin{itemize}
    \item[\textbf{R-F-1}] \textbf{Natural language querying:} the system accepts a
        free-text English question and returns a textual answer with source citations.
        When the answer is absent from the indexed documents, the system must say so
        explicitly rather than generating an unsupported response.

    \item[\textbf{R-F-2}] \textbf{Model selection:} the user may select which language
        model handles the current query. The system must apply the selected model to all
        subsequent queries in the session without requiring a restart.

    \item[\textbf{R-F-3}] \textbf{Document upload:} the user may upload a document
        that is immediately chunked, embedded, and added to the vector store, making
        its contents available for retrieval in the same session.

    \item[\textbf{R-F-4}] \textbf{Embedding and retrieval:} each chunk is encoded into
        a dense vector and stored in a local persistent vector database. At query time
        the top-$K$ most similar chunks are retrieved; $K$ is configurable.

    \item[\textbf{R-F-5}] \textbf{Session history:} the system persists conversation
        history across page reloads. The user may view past sessions, continue a
        previous session, or delete a session entirely.

    \item[\textbf{R-F-6}] \textbf{System information:} the user may view the current
        system status, including model availability and vector store state, without
        submitting a query.

    \item[\textbf{R-F-7}] \textbf{REST API:} the backend exposes a \texttt{/health}
        endpoint reporting system status and a \texttt{/chat} endpoint that accepts a
        message and returns the answer with source references. Both endpoints must be
        usable independently of the browser interface.

    \item[\textbf{R-F-8}] \textbf{Web interface:} a browser-based chat UI displays
        conversation history, source file names, model availability, and inline error
        messages without requiring any installation on the client side.
\end{itemize}

% ─────────────────────────────────────────────────────────────
\section{Non-Functional Requirements}
\label{sec:req-nonfunctional}

\begin{itemize}
    \item[\textbf{R-NF-1}] \textbf{Privacy:} no query, chunk, or answer may leave the
        local machine. The system must remain fully functional in a network-isolated
        environment after initial setup.

    \item[\textbf{R-NF-2}] \textbf{Resource constraints:} the system must run on 8~GB
        RAM with no discrete GPU, using a quantised (4-bit or 8-bit) language model and
        a CPU-friendly embedding model.

    \item[\textbf{R-NF-3}] \textbf{Response latency:} end-to-end response time must
        not exceed 90~seconds on the reference hardware. Embedding and retrieval must
        complete in under one second; generation typically takes 30--80~seconds on CPU.

    \item[\textbf{R-NF-4}] \textbf{Offline operation:} after the initial model pull
        and corpus crawl, no network connection is required. All subsequent pipeline
        steps run entirely from local storage.

    \item[\textbf{R-NF-5}] \textbf{Reliability:} if the language model process is
        unavailable, \texttt{/health} must report it and \texttt{/chat} must return an
        HTTP~503. The vector store and log database must not be corrupted on unexpected
        shutdown, and re-running the pipeline must be idempotent.

    \item[\textbf{R-NF-6}] \textbf{Configurability:} all tunable parameters — model
        names, $K$, temperature, paths, timeouts — are defined in a single
        configuration module. Swapping models requires only a one-line change there;
        changing the embedding model additionally requires rebuilding the vector store.

    \item[\textbf{R-NF-7}] \textbf{Deployability:} the system must be deployable on
        a faculty server without modification to the application code. All components
        are containerised with Docker, so that moving from a local prototype to a live
        university service requires infrastructure provisioning only — configuring
        ports, paths, and model names — with no changes to any other part of the
        codebase.
\end{itemize}